% 2013 四川卷文科第 18 题：输出 y 的值的程序框图。
\begin{tikzpicture}[font=\footnotesize, >=Stealth, line width=0.45pt,
  node distance=4.6mm,
  term/.style={draw, rounded corners=5pt, minimum height=1.8em, inner xsep=8pt},
  io/.style={draw, trapezium, trapezium left angle=72, trapezium right angle=108,
             minimum height=1.8em, inner xsep=7pt},
  box/.style={draw, minimum height=1.8em, inner xsep=6pt},
  dec/.style={draw, diamond, aspect=2.4, inner sep=1pt}]
  \node[term] (start) {开始};
  \node[io, below=5mm of start] (input) {输入 $x$};
  \node[dec, below=6mm of input] (even) {$x$ 为偶数?};
  \node[box, right=9mm of even] (y1) {$y=1$};
  \node[dec, below=9mm of even] (div) {$x$ 能被 3 整除?};
  \node[box, left=9mm of div] (y2) {$y=2$};
  \node[box, below=6mm of div] (y3) {$y=3$};
  \node[io, below=15mm of y3] (output) {输出 $y$};
  \node[term, below=5mm of output] (end) {结束};
  \draw[->] (start) -- (input);
  \draw[->] (input) -- (even);
  \draw[->] (even) -- node[above] {否} (y1);
  \draw[->] (even) -- node[right] {是} (div);
  \draw[->] (div) -- node[above] {否} (y2);
  \draw[->] (div) -- node[right] {是} (y3);
  \draw[->] (y1.south) |- (output.east);
  \draw[->] (y2.south) |- (output.west);
  \draw[->] (y3) -- (output);
  \draw[->] (output) -- (end);
\end{tikzpicture}
